\section{Do Not Become Another Person's Tempter}

Christ places the warning about the little ones immediately before the warning about the hand, foot, and eye. He speaks with terrible severity of the person who causes one of his believing little ones to fall (Mk 9:41). Renunciation must therefore include conduct by which we endanger someone else. A person may be concerned about resisting his own temptations and still persuade another to do what he knows is wrong. The other's salvation also places a claim upon his conduct. Pleasure, convenience, or standing among others cannot justify drawing a soul away from God.

The Church calls this evil scandal. The word here concerns leading someone into sin. Public embarrassment or the discovery of a fault may accompany it, but its deepest harm is spiritual. A person can cause scandal quietly, in a conversation no one else hears, by making evil attractive or by teaching another to dismiss the voice of conscience. The Catechism describes the scandal-giver as the neighbor's tempter and teaches the gravity of deliberately leading another into a grave offense (CCC 2284).

St.~Thomas considers scandal among the sins opposed to charity. He begins with the image of an obstacle placed in a traveler's path: a word or deed can become a stumbling block on the spiritual way. Commands, inducements, and examples can all draw another toward sin (\emph{Summa theologiae} II--II, q.~43, a.~1). Each possibility asks something different of the conscience. We may have ordered what was wrong, persuaded someone to join us, or made our own conduct appear a safe pattern for imitation.

Aquinas also preserves the distinction between our influence and another person's will. Our word can provide an occasion for his sin; it does not replace his consent (a.~1, reply~3). This keeps responsibility truthful on both sides. We cannot excuse our deliberate encouragement by saying that the other person chose freely. Neither should we imagine that every later fault in someone we have known was secretly caused by us. The examination concerns the real help we gave to evil, including what our position or example reasonably communicated.

Authority makes this responsibility more searching. A child receives more than instructions from a parent; the child also observes which instructions the parent treats as binding. A younger worker learns what integrity costs by seeing whether a superior protects it under pressure. The Catechism therefore attends to the authority of the scandal-giver, the weakness of the person influenced, and the special duties of those who teach or educate (CCC 2285--2287). Using such trust to undermine a conscience betrays the person who relies upon us.

The temptation can be concealed beneath an offer of belonging. A friend hesitates to join a dishonest practice, and we work upon his hesitation until he is ashamed of resisting. We may tell ourselves that we are relieving an unnecessary scruple, though we know why the act is wrong. His refusal also exposes our own conduct, and persuading him would make that exposure easier to dismiss. The invitation thus seeks two advantages for us at his expense: companionship in the act and relief from its accusation. Concern for his salvation requires us to respect the conscience we were trying to silence.

The influence of example can outlast the act. An excuse taught in a household may be carried into another household by someone who learned it there. We cannot calculate every later consequence, but we can examine the conduct we taught people to admire or dismiss. A younger person may have learned that fidelity is foolish because we regularly mocked it. Repentance then requires attention to the lesson as well as the original act. If our example made resistance harder, a candid acknowledgment of its wrongness may help the person we influenced recover confidence in what his conscience told him.

Care for another's soul also requires discernment. Aquinas distinguishes giving scandal from another person's taking occasion to sin even from a good act (a.~1, reply~4). Someone may resent a just refusal or envy another's virtue. That reaction does not make obedience an injury. We are not required to surrender God's command whenever it displeases someone. Otherwise concern for scandal could itself become a way of compelling sin, with the most unreasonable person determining the conduct of everyone nearby.

In treating spiritual goods, Thomas refuses the proposal that one should commit sin in order to prevent someone else's sin. He also asks for patience with the weak. An action which can prudently be deferred, or a matter which needs explanation, should be handled with attention to the person's capacity and need. The truth must be taught faithfully, with judgment concerning persons and occasions (a.~7). This leaves room for gentleness, a quieter manner, and a better time. It leaves no room for using charity as the name of a lie or for presenting an evil as good.

Repentance for scandal accordingly seeks the neighbor's good as well as our own pardon. A harmful instruction may need to be withdrawn. A person whom we pressured may need to hear that his original hesitation was right. An arrangement which rewards dishonesty may need to be changed. Such correction is to be proportionate and prudent; it need not spread a hidden sin to people who had no knowledge of it. The aim is to stop providing help to evil and, where possible, remove the false lesson our conduct has given.

A good example can also assist another person's freedom. When a superior accepts a loss rather than demand wrongdoing, those beneath him see that integrity can survive its cost. He removes the pressure his authority might have added to temptation. A parent who admits a fault likewise makes it easier for a child to acknowledge the truth about his own. Christ promises a reward for the cup of water given because someone belongs to him. Such assistance serves a person whom he calls to eternal life, and belongs beside the refusal to obstruct that person's way to God.

